\documentclass[10pt,a4paper]{amsart}
\usepackage[margin=19mm]{geometry}
\usepackage{amsmath,amssymb,mathtools}
\usepackage[scr=rsfs,cal=euler]{mathalfa}
\usepackage{xfrac}
\usepackage[unicode,hidelinks,bookmarks=false]{hyperref}
\newcommand{\ie}{i.e.\ }
\newcommand{\cf}{cf.\ }
\newcommand{\resp}{respectively\ }
\newcommand{\Subsection}{\subsection}
\providecommand{\nobreakdash}{\nobreak\hbox{-}}
\setlength{\emergencystretch}{2em}
\multlinegap0pt
\usepackage{bm}
\usepackage{bbm}
\usepackage{dsfont}
\usepackage{xspace}
\usepackage{cancel}
\usepackage{mathtools}
\usepackage{amsthm}
\makeatletter
\@ifundefined{Lemma}{\newtheorem{Lemma}{Lemma}}{}
\@ifundefined{Proposition}{\newtheorem{Proposition}{Proposition}}{}
\makeatother
\newcommand{\E}{\mathbb{E}}
\newcommand{\PP}{\mathbb{P}}
\newcommand{\R}{\mathbb{R}}
\newcommand{\abs}[1]{\left\lvert#1\right\rvert}
\newcommand{\bbI}{\mathbb{I}}
\newcommand{\cD}{\mathcal{D}}
\newcommand{\cH}{\mathcal{H}}
\title[Multiple Testing of Linear Forms for Noisy Matrix Completion]{Multiple Testing of Linear Forms for Noisy Matrix Completion}
\author{Wanteng Ma, Lilun Du, Dong Xia, Ming Yuan}
\date{Source: arXiv:2312.00305v3 (CC BY 4.0)}
\begin{document}
\maketitle

\noindent\emph{Source and licence.} Multiple Testing of Linear Forms for Noisy Matrix Completion. Version arXiv:2312.00305v3 (CC BY 4.0), distributed under the Creative Commons Attribution 4.0 International licence, as recorded in the arXiv licence metadata for this exact version and shown on the abstract page https://arxiv.org/abs/2312.00305v3. Full rights notice: licenses/arxiv-2312.00305-CC-BY-4.0.txt.\par\smallskip

\noindent\textit{Editorial scope.} Counted unit: the Lemma whose source label is \texttt{lemma:conv-prob} in the article (source0.tex lines 1978--1983, source label \texttt{lemma:conv-prob}), delivered together with the article's own complete original proof of it (source0.tex lines 1985--2029, one \texttt{proof} environment of 45 lines). No mathematics is written, repaired or re-derived: statement and proof are the article's own text line for line. Required context retained uncounted: eq:weak-corr-1 (lines 532--536); prop:asy-w1 (lines 1723--1731); lemma:weak-cov (lines 1790--1798). Every definition, display and label of those lines is retained unchanged. Omitted from this package and not redistributed: 1--531, 537--1722, 1732--1789, 1799--1977, 1984--1984, 2030--3070. Nothing inside a retained excerpt is omitted or altered: every retained body is delivered exactly as the article's lines are written, so no omission receipt of any kind accompanies this package. Citations: the retained excerpts carry no citation command at all, so no bibliography entry of the article is transcribed and no reference is redistributed. Reference adaptation: none was needed and none was made; every reference inside the delivered unit resolves inside it, so no unresolved reference or citation is produced. Printed slips kept literal: no printed word, symbol, hypothesis or number of the counted statement or of its proof was altered. Layout and class: the article's own class is \texttt{article} with packages including geometry, inputenc, amsfonts, amsgen, amsmath, amstext, amsbsy, amsopn, amssymb, subcaption, graphicx, comment, enumitem, natbib; the wrapper is the lane's pinned \texttt{amsart} editorial wrapper at 10pt on A4, which restates the theorem-like environments the retained text needs and adds the article's own macro definitions that the retained text needs (\texttt{\textbackslash E}, \texttt{\textbackslash PP}, \texttt{\textbackslash R}, \texttt{\textbackslash abs}, \texttt{\textbackslash bbI}, \texttt{\textbackslash cD}, \texttt{\textbackslash cH}), with extra wrapper packages (bm, bbm, dsfont, xspace, cancel, mathtools). The delivered numbering is the unit's own. Assets: no retained span contains a figure environment, a graphics-inclusion command, a PDF-page inclusion, a table or a TikZ picture; no image file is copied into this package, the package declares no asset dependency and no image file is redistributed with it. Licence: version arXiv:2312.00305v3 of this article is distributed under the Creative Commons Attribution 4.0 International licence, as recorded in the arXiv licence metadata for this exact version and shown on the abstract page https://arxiv.org/abs/2312.00305v3. A full rights notice is delivered at licenses/arxiv-2312.00305-CC-BY-4.0.txt. Characters: every retained excerpt is pure ASCII, so the delivered engine, XeLaTeX with its default Latin Modern text font, reports no missing character.
\begin{equation}\label{eq:weak-corr-1}
	\begin{aligned}
		\cH_{0,\text{strong} }^2 := \left\{ (T_1,T_2)\in\cH_0\times \cH_0 : \rho_{T_1,T_2}\ge c q_0^{-\nu} \right\},
	\end{aligned}
\end{equation}

\begin{Proposition}\label{prop:asy-w1}
There exits a constant $C_2$ such that $W^{(1)}_T$ follows the asymptotic normality rate:
\begin{equation}\label{eq:asy-w1}
    \begin{aligned}
    \sup_{t\in\R}\abs{\PP(W^{(1)}_T>t )-\Phi(-t)}\le C_2 h_n,
    \end{aligned}
\end{equation}
for any $T\in \cH_0$.
\end{Proposition}

\begin{Lemma}[Weak dependency of null statistics]\label{lemma:weak-cov}
        Conditional on $\cD_1$, 
\begin{equation}
\begin{aligned}
          \sup_{0\le t\le L_n } \frac{\sum_{(T_i,T_j)\in \cH_{0,\text{weak} }^2 } \abs{\operatorname{cov}(\bbI(W_{T_i}^{\mathsf{rank} }>t),\bbI(W_{T_j}^{\mathsf{rank} }>t)) }}{ q_0^2 G^2(t)}\le  C_1 \frac{h_n q_0}{\epsilon_n \eta_n }+C_2 \frac{1}{ \left(\epsilon_n\eta_n q_0\right)^{\nu/2} },
\end{aligned}
\end{equation}
where $\nu$ is the weak correlation parameter defined in \eqref{eq:weak-corr-1}.
\end{Lemma}

\begin{Lemma}\label{lemma:conv-prob}
For any $\varepsilon>0$, conditional on $\cD_1$, it holds that
\begin{equation*}
   \PP\left(\sup\limits_{0\le t\le L_n}\abs{\frac{ \sum_{T\in\cH_0}\bbI(W_T^{\mathsf{rank} }>t ) }{q_0 G(t) } -1} \ge \varepsilon\right) \le \frac{C}{\varepsilon^2} \log(\frac{q_0 }{\epsilon_n \eta_n}) \left( \left(\frac{\beta_{\mathsf{s}} q_0^2 }{ \epsilon_n^2\eta_n^2}\right)^{\frac{1}{2}} + \left(\frac{h_n q_0}{\epsilon_n \eta_n}+ \frac{1}{ \left(\epsilon_n\eta_n q_0\right)^{v/2} } \right)^{\frac{1}{2}}\right) .
\end{equation*}
\end{Lemma}

\begin{proof}
To prove the uniform convergence in probability, we define a grid on $[0,L_n]$:
%with the same spirit as \cite{du2021false}: 
$$\left\{ t_k= G^{-1}\left( \frac{1}{2} (2G(L_n))^{\frac{k}{K}} \right) \right\}_{k=0}^{K},  $$
which equates each $G(t_k)$ with $\frac{1}{2} (2G(L_n))^{\frac{k}{K}}$. Then for each $t\in [t_{k-1},t_{k})$, the ratio can be bounded by:
\begin{equation*}
    \frac{ \sum_{T\in\cH_0}\bbI(W_T^{\mathsf{rank} }>t_{k} ) }{q_0 G(t_{k-1}) } \le \frac{ \sum_{T\in\cH_0}\bbI(W_T^{\mathsf{rank} }>t ) }{q_0 G(t) } \le  \frac{ \sum_{T\in\cH_0}\bbI(W_T^{\mathsf{rank} }>t_{k-1} ) }{q_0 G(t_{k}) }.
\end{equation*}
Define $(2G(L_n))^{\frac{1}{K}}=r_K$, we have $G(t_k)/G(t_{k-1})=r_K$, and 
\begin{equation*}
    \abs{\frac{ \sum_{T\in\cH_0}\bbI(W_T^{\mathsf{rank} }>t ) }{q_0 G(t) } -1}\le \sup_{i=k-1,k}\frac{1}{r_K} \abs{\frac{ \sum_{T\in\cH_0}\bbI(W_T^{\mathsf{rank} }>t_i ) }{q_0 G(t_{i}) }-1} +\abs{r_K-1}\vee\abs{\frac{1}{r_K}-1 },
\end{equation*}
for each $t\in [t_{k-1},t_{k})$. Then for any $t\in [0,L_n]$, it suffices to control the quantities 
\begin{equation*}
    \sup_{k=0,\dots, K}\frac{1}{r_K} \abs{\frac{ \sum_{T\in\cH_0}\bbI(W_T^{\mathsf{rank} }>t_k ) }{q_0 G(t_{k}) }-1}\le  \sup_{k=0,\dots, K}\frac{1}{r_K} \abs{\frac{ \sum_{T\in\cH_0}\bbI(W_T^{\mathsf{rank} }>t_k )-\PP(W_T^{\mathsf{rank} }>t_k ) }{q_0 G(t_{k}) }} + C_3 \frac{h_n q_0}{\epsilon_n \eta_n }
\end{equation*}
and $\abs{r_K-1}\vee\abs{\frac{1}{r_K}-1 }$ by Proposition \ref{prop:asy-w1}.  Denote 
$$D_n=\sup\limits_{k=0,\dots, K}\frac{1}{r_K} \abs{\frac{ \sum_{T\in\cH_0}\bbI(W_T^{\mathsf{rank} }>t_k )-\PP(W_T^{\mathsf{rank} }>t_k ) }{q_0 G(t_{k}) }}.$$ 
It follows that
\begin{equation}
    \begin{aligned}
              &\E D_n^2\le \frac{K}{r_K^2}\E \abs{\frac{ \sum_{T\in\cH_0}\bbI(W_T^{\mathsf{rank} }>t_{k} )-\PP(W_T^{\mathsf{rank} }>t_{k} ) }{q_0 G(t_{k}) }}^2\\
              &\le  \frac{K}{r_K^2} \frac{\sum\limits_{(T_1,T_2)\in\cH^2_{0,\text{weak}}} \abs{\operatorname{cov}(\bbI(W_{T_1}^{\mathsf{rank} }>t),\bbI(W_{T_2}^{\mathsf{rank} }>t)) } +\sum\limits_{T_1,T_2\in\cH^2_{0,\text{strong}} } \abs{\operatorname{cov}(\bbI(W_{T_1}^{\mathsf{rank} }>t),\bbI(W_{T_2}^{\mathsf{rank} }>t)) } }{q_0^2 G^2(t) },
    \end{aligned}
\end{equation}
for any $t\in \{t_k\}$. Since the number of strong dependency pairs $\abs{\cH^2_{0,\text{strong}}}\le \beta_{\mathsf{s}}q_0^2 $, we have
\begin{equation*}
    \frac{\sum\limits_{T_1,T_2\in\cH^2_{0,\text{strong}} } \abs{\operatorname{cov}(\bbI(W_{T_1}^{\mathsf{rank} }>t),\bbI(W_{T_2}^{\mathsf{rank} }>t)) }}{q_0^2 G^2(t) }\le \frac{\beta_{\mathsf{s}}q_0^2 }{ \epsilon_n^2\eta_n^2},
\end{equation*}
for any $t\in [0,L_n]$. For the weak dependency pair, 

\begin{equation*}
    \frac{\sum\limits_{T_1,T_2\in\cH^2_{0,\text{weak}} } \abs{\operatorname{cov}(\bbI(W_{T_1}^{\mathsf{rank} }>t),\bbI(W_{T_2}^{\mathsf{rank} }>t)) }}{q_0^2 G^2(t)}\le C_1 \frac{h_n q_0}{\epsilon_n \eta_n }+C_2 \frac{1}{ \left(\eta_n q_0\right)^{v/2} },
\end{equation*}
 where we apply our previous results in Lemma \ref{lemma:weak-cov}. What remains for us is to specify the density of grid $K$. Choose a constant $\varsigma$ and we set
 $$K= \log(\frac{q_0 }{\epsilon_n \eta_n}) \min \left\{ \left(\frac{ q_0^2 \beta_{\mathsf{s}} }{ \eta_n^2 \epsilon_n  }\right)^{-\varsigma} , \left( \frac{q_0 h_n}{\eta_n\epsilon_n}+\frac{1}{ \left(\epsilon_n\eta_n q_0\right)^{v/2} } \right)^{-\varsigma} \right\}  ,$$ 
 then it is clear that $1\le\frac{1}{r_k}\le \left[ \frac{q_0}{\epsilon_n \eta_n} \right]^{1/K}\to 1$, and $ K (\frac{\beta_{\mathsf{s}}q_0^2 }{ \epsilon_n^2\eta_n^2} + \frac{h_n q_0}{\epsilon_n \eta_n })\to 0$. Therefore 
 \begin{equation*}
\begin{aligned}
      &\abs{r_K-1}\vee\abs{\frac{1}{r_K}-1 } \le C \frac{1}{K}\log(\frac{q_0 }{\epsilon_n \eta_n}) \le  \left( \left(\frac{\beta_{\mathsf{s}} q_0^2 }{ \epsilon_n^2\eta_n^2}\right)^{\varsigma} + \left(\frac{h_n q_0}{\epsilon_n \eta_n} +\frac{1}{ \left(\epsilon_n\eta_n q_0\right)^{v/2} }\right)^{\varsigma}\right) \\
    &\E D^2_n  \le C K \left(\frac{\beta_{\mathsf{s}}q_0^2 }{ \eta_n^2} + \frac{h_n q_0}{\epsilon_n \eta_n} \right)\le C \log(\frac{q_0 }{\epsilon_n \eta_n}) \left( \left(\frac{\beta_{\mathsf{s}} q_0^2 }{ \epsilon_n^2\eta_n^2}+\frac{1}{ \left(\epsilon_n\eta_n q_0\right)^{v/2} }\right)^{1-\varsigma} + \left(\frac{h_n q_0}{\epsilon_n \eta_n} \right)^{1-\varsigma}\right).
\end{aligned}
 \end{equation*}
We can finish the proof of uniform convergence by using Markov's inequality with $\varsigma=\frac{1}{2}$.
\end{proof}

\end{document}
